% ======================================================================
%  2. Background and critical state of the art
% ======================================================================
\section{Background and critical state of the art}
\label{sec:background}

This section introduces the material required
to understand the contributions, and organizes the literature by
\emph{scientific question} rather than paper by paper. The theoretical thread we
follow is a single argument,
\[
\begin{gathered}
\underbrace{\text{well-specified GP-UCB}}_{\S\ref{sec:bg-gpucb}}
\;\rightarrow\;
\underbrace{\text{composite BO}}_{\S\ref{sec:bg-composite}}
\;\rightarrow\;
\underbrace{\eps\text{-misspecified GP bandits}}_{\S\ref{sec:bg-misspec}}
\\[0.1cm]
\;\rightarrow\;
\underbrace{\text{convergence despite smoothness misspecification}}_{\S\ref{sec:bg-misspec}}
\\[0.1cm]
\;\rightarrow\;
\boxed{\epsinf(B)\text{ barrier (Sec.~\ref{sec:contrib-tm})}},
\end{gathered}
\]
\vspace{0.1cm}
which makes precise both why the initial regularity obstruction is only a first
step and why the quantitative approximation barrier is the stronger statement it
led to. It then reviews papers to provide guidelines on the specific problems that we study (multiobjective, alpha-prior BO).

\subsection{Gaussian-process Bayesian optimization}
\label{sec:bg-gpucb}

We consider an unknown objective $f:\X\to\R$ on a compact domain
$\X\subset\R^d$, observed through noisy queries $y_t=f(x_t)+\eps_t$. A Gaussian
process prior $f\sim\GP(0,k)$ yields, after $t-1$ observations, a posterior with
mean $\mu_{t-1}$ and standard deviation $\sigma_{t-1}$ available in closed form.
The GP-UCB rule of Srinivas et al.~\cite{srinivas2010} selects
\begin{equation}
  x_t \;\in\; \argmax_{x\in\X}\;\bigl(\mu_{t-1}(x)+\sqrt{\beta_t}\,\sigma_{t-1}(x)\bigr),
  \label{eq:bg-gpucb}
\end{equation}
trading exploitation ($\mu_{t-1}$) against exploration ($\sigma_{t-1}$). The
statistical difficulty of a kernel is captured by the \emph{maximum information
gain} $\gain_T=\max_{|A|=T}\tfrac12\log\det(I+\sigma^{-2}K_A)$, and the central
guarantee is that, with high probability, cumulative regret is
$R_T=O(\sqrt{T\beta_T\gain_T})$
\cite{srinivas2010}. \hfill\Plit

We recall the ingredients (posterior formulas, information gain, the
high-probability confidence bound and the regret theorem) in
Appendix~\ref{app:background}; for the GP/RKHS/kernel-ridge connections we follow
the review of Kanagawa et al.~\cite{kanagawa2018}. \hfill\Plit

\subsection{RKHS assumptions and why kernel choice is a regularity assumption}
\label{sec:bg-rkhs}

The regret guarantee rests on an assumption that is easy to state and easy to
overlook: the objective belongs to the reproducing-kernel Hilbert space (RKHS)
$\Hk$ of the kernel, with bounded norm, $f\in\Hk$ and $\Hnorm{f}\le B$. This is
not a technicality. Choosing a kernel \emph{is} choosing a function space, and
hence a regularity assumption on $f$. For stationary kernels the RKHS is a
Sobolev-type space whose smoothness is dictated by the decay of the spectral
density: the Mat\'ern-$\nu$ kernel gives $\Hk\simeq H^{\nu+d/2}(\R^d)$, while the
RBF kernel gives functions that are $C^\infty$ (indeed real-analytic)
\cite{kanagawa2018}. A classical regularity-transfer result makes the link
between kernel smoothness and RKHS smoothness explicit.

\begin{theorem}[Regularity transfer, {see Appendix~\ref{app:background}}]
\label{thm:bg-regtransfer}
If $k\in C^{2m}(\X\times\X)$ then $\Hk\subset C^m(\X)$; every $f\in\Hk$ is
$C^m$, and $\partial^\alpha f(x)=\inner{f}{\partial_x^\alpha k(x,\cdot)}_{\Hk}$
for $|\alpha|\le m$.
\provbadge{colSpec}{literature, specialized here}
\end{theorem}

\noindent Revisiting this transfer result was our first route toward
understanding task-induced misspecification: if the task transformation destroys
differentiability, and the kernel forces $C^m$ regularity, the transformed
objective cannot lie in $\Hk$. We treat Theorem~\ref{thm:bg-regtransfer} as
(specialized) background, not as a novel result, and use it only to motivate the
target-matching analysis of Section~\ref{sec:contrib-tm}.

\subsection{Composite Bayesian optimization: a direct predecessor}
\label{sec:bg-composite}

The idea of modelling a latent quantity rather than the objective actually
optimized is \emph{not} new to this work. Astudillo and
Frazier~\cite{astudillo2019} study \emph{composite} objectives $g(h(x))$, where
$h$ is expensive and $g$ is known and cheap; they model $h$ with a GP and apply
$g$ to the induced posterior when forming the acquisition, and they prove
asymptotic consistency for their EI-CF algorithm. Their examples include
calibration objectives $\norm{h(x)-y_{\mathrm{obs}}}$, which are structurally very
close to target matching. \hfill\Plit

We therefore do \emph{not} claim that ``model the phenomenon rather than the
composite objective'' is a novel idea. The gap we address is different: composite
BO exploits \emph{known} composite structure for statistical efficiency, but does
not characterize the RKHS regularity mismatch created by directly modelling a
\emph{non-smooth} task transformation, nor the norm--accuracy trade-off such a
transformation induces. Our target-matching analysis investigates precisely this
modelling distinction, through RKHS misspecification, a quantitative
approximation barrier, and confidence-set pushforward pushforward arguments that preserve valid uncertainty bounds and yield convergence guarantees for UCB-based acquisition policies.
\hfill\Psynth

\subsection{Misspecification: leaving the RKHS need not mean failure}
\label{sec:bg-misspec}

Suppose the optimized function falls outside the assumed RKHS. Two very different
things can then be true, and they must be kept apart:
\begin{enumerate}
  \item a standard \emph{guarantee} is lost (the confidence bound and regret
        proof of \S\ref{sec:bg-gpucb} no longer apply);
  \item the algorithm \emph{actually fails} (e.g.\ incurs linear regret).
\end{enumerate}
Two lines of work delineate this gap and are the direct bridge to our
contribution.

\paragraph{$\eps$-misspecified GP bandits.}
Bogunovic and Krause~\cite{bogunovic2021} consider a true function $g$ that is
only \emph{approximately} in the RKHS,
$\inf_{\Hnorm{h}\le B}\norm{g-h}_\infty\le\eps$, and derive robust GP-bandit
guarantees whose regret degrades gracefully in the misspecification level $\eps$.
This reframes the right question: the relevant quantity is not whether
$g\in\Hk$, but \emph{how accurately} $g$ can be approximated by an RKHS ball of
controlled radius. \hfill\Plit

\paragraph{Convergence despite smoothness misspecification.}
Wynne, Briol and Girolami~\cite{wynne2021} show, complementarily, that
\[
g\notin\Hk \;\;\not\Rightarrow\;\; \text{GP approximation cannot converge}.
\]
They give convergence rates for GP posterior means under smoothness
misspecification, controlled by the geometry of the design (fill distance,
quasi-uniformity), and derive BO implications under exploration strategies that
produce sufficiently regular designs. Their analysis can even accommodate
sequences of adaptively chosen kernel hyperparameters, under explicit bounds on
the resulting smoothness and norm-equivalence constants --- which should
\emph{not} be read as a licence for arbitrary marginal-likelihood tuning.
\hfill\Plit

\paragraph{Consequence for our programme.}
These two results are exactly why the report does not stop at a qualitative
non-membership theorem. They motivate the progression
\emph{loss of differentiability} $\to$ \emph{RKHS non-membership} $\to$
\emph{recognition that out-of-RKHS functions may still be consistently
approximated} $\to$ \emph{a quantitative, norm-constrained approximation
barrier}. The barrier (Section~\ref{sec:contrib-tm}) answers the inverse of
Bogunovic--Krause's question for the structured transformation $g=\abs{f-T}$:
how rapidly must RKHS complexity grow to approximate the task-transformed
objective? We are careful \emph{not} to combine Bogunovic--Krause's generic
misspecification lower bounds mechanically with our barrier to claim a lower
regret bound for task-in-model GP-UCB: their lower bounds concern a broader
adversarial class and do not automatically specialize to the family
$g=\abs{f-T}$. \hfill\Psynth

\paragraph{Online hyperparameter estimation.}
Our experiments re-estimate GP hyperparameters online by marginal likelihood. 
However, the standard GP-UCB guarantees used in our theoretical analysis assume
fixed, known kernel parameters together with an appropriate RKHS-norm bound.
Online MLE therefore falls outside the scope of these guarantees in general.
Berkenkamp et al.~\cite{berkenkamp2019} show that naive hyperparameter adaptation
can lead to over-smooth models and insufficient exploration, and recover
guarantees by deliberately expanding the candidate function class over time.
Bull's analysis of expected improvement~\cite{bull2011} similarly illustrates
how data-driven scale estimation can suppress exploration. We use these results
to contextualize the behavior observed in our experiments, rather than as part
of the main theoretical guarantees developed here.\hfill\Plit

\subsection{Multi-objective Bayesian optimization}
\label{sec:bg-mo}

With a vector objective $F=(f_1,\dots,f_m)$ there is no single optimum but a
\emph{Pareto front} of non-dominated trade-offs. The standard reduction is
\emph{scalarization}: pick a family $\{\prm_\lambda:\R^m\to\R\}$ and solve
$\max_x\prm_\lambda(F(x))$ with any single-objective optimizer. Two properties
matter. A scalarization is \emph{sound} if every maximizer is Pareto-optimal
(guaranteed by coordinate-wise monotonicity), and \emph{complete} if varying
$\lambda$ traces out the whole front. Linear scalarizations are sound but not
complete --- they only reach the convex hull of the front --- whereas
Chebyshev and hypervolume scalarizations are complete. \hfill\Plit

For random-scalarization methods we build on Paria et al.~\cite{paria2019}
(random scalarizations for multi-objective BO) and, centrally, on Golovin and
Zhang~\cite{golovin2020}. Golovin and Zhang introduce \emph{random hypervolume
scalarizations} and prove a remarkable identity rewriting the dominated
hypervolume as an expectation of a maximum of scalarized values; from it they
obtain UCB/Thompson-sampling convergence guarantees and a generic reduction from
a convergent single-objective optimizer to Pareto-front approximation.
\hfill\Plit

\begin{quote}
\itshape Attribution.\ The random hypervolume scalarization, the
regret--hypervolume identity, the associated convergence guarantees, and the
generic single-objective-to-Pareto reduction are due to Golovin and
Zhang~\cite{golovin2020} and are used here as imported theory. Our contribution
in the multi-objective setting is the separation-of-concerns interpretation and
the pushforward-confidence analysis of these constructions
(Section~\ref{sec:contrib-mo}), not the invention of the scalarization machinery.
\end{quote}

Differentiable expected-hypervolume-improvement methods
\cite{daulton2020} are an alternative, direct route to modern multi-objective
BO; they are peripheral to our modelling question and are mentioned only for
completeness. The hypervolume indicator itself, and the strength-Pareto
viewpoint, go back to Zitzler and Thiele~\cite{zitzler1999}. \hfill\Plit

\subsection{Prior-guided Bayesian optimization: two meanings of ``prior''}
\label{sec:bg-prior}

Prior knowledge enters BO in two fundamentally different ways, and conflating
them is exactly the failure mode we study.

\paragraph{A spatial prior at the decision level.}
Guay Hottin et al.~\cite{guayhottin2025} encode prior information about the
probable \emph{location} of the optimizer through an acquisition modulation of
the $\alpha$-$\pi$BO form,
\[
A_t(x)\,\bigl[\alpha\,\pi(x)+(1-\alpha)\bigr],
\]
building on the user-belief acquisition $\pi$BO of Hvarfner et
al.~\cite{hvarfner2022}. For expected improvement the method retains the standard
convergence order up to a controlled factor $1/(1-\alpha)$. This is the canonical
example of a \emph{spatial} prior acting at the decision level: because the
weight is bounded below, even a misspecified prior cannot eliminate any region of
the search space. \hfill\Plit

\paragraph{A functional prior at the model level.}
Cheng et al.~\cite{cheng2024} instead assume that a complete surrogate function
$f'$ is informative about the \emph{entire} objective and model
$f\sim\GP(\lambda f',k)$. For fixed $\lambda$ the relevant complexity becomes
$\norm{f-\lambda f'}_{\Hk}$, minimized at the oracle RKHS-projection coefficient
$\lambda^\star_{\rkhs}=\inner{f}{f'}_{\Hk}/\norm{f'}^2_{\Hk}$. This is important
because it \emph{prevents} an overly broad claim that non-zero GP means or
function priors are intrinsically undesirable: a genuine global functional prior
can be theoretically beneficial. \hfill\Plit

\paragraph{The gap we address.}
Our critique is narrower than ``model-level priors are bad''. It is that
\begin{quote}
\itshape information about where the optimizer is likely to lie is not, without
additional assumptions, information about the entire response surface.
\end{quote}
We use Cheng et al.\ to motivate an analysis of the gap between the inaccessible
\emph{global} RKHS criterion and coefficients estimated only from adaptively
sampled observations, and we clearly distinguish their proposed estimator from
our subsequent local-versus-global and self-confirmation analysis
(Section~\ref{sec:contrib-prior}). \hfill\Psynth

\subsection{Critical synthesis: the gap addressed by this internship}
\label{sec:bg-synthesis}

The results above each study a \emph{particular} downstream task or a
\emph{generic} misspecification class. Composite BO exploits known structure for
efficiency; $\eps$-misspecified and smoothness-misspecified analyses tell us that
leaving the RKHS costs a guarantee but need not cost convergence; random
scalarizations turn single-objective optimizers into Pareto solvers; spatial and
functional priors sit at opposite ends of the ``where does the prior act''
spectrum. What is missing is a study of the \emph{modelling boundary} itself ---
the line between the latent phenomenon $f$ and the task-specific decision
transformation $\prm$ --- viewed uniformly across tasks. Placing that boundary is
the object of the next three sections.
